\subsubsection{内点係数の Taylor 展開による導出}
\label{app:ccd_coef_derivation}
% ═══════════════════════════════════════════════════════════════════════════════

本付録は §\ref{sec:ccd_te_I} で述べた CCD スキームの係数
$(\alpha_1, a_1, b_1)$ および $(\beta_2, a_2, b_2)$ を，
Taylor 展開による打ち切り誤差解析から厳密に導出する．
各 Equation の RHS を格子点 $i$ 周りで展開し，
LHS との係数比較により3元連立方程式を立て，代数的に解く．

\subsubsection{Equation-I の係数導出}
\label{app:ccd_coef_derivation_I}

Equation-I（式~\eqref{eq:CCD_I}）の各項を $f$ の Taylor 展開で表す．
添字 $i\pm1$ の一次微分値と二次微分値は：
\begin{align}
  f'_{i\pm1} &= f'_i \pm hf''_i + \frac{h^2}{2}f^{(3)}_i \pm \frac{h^3}{6}f^{(4)}_i
               + \frac{h^4}{24}f^{(5)}_i \pm \frac{h^5}{120}f^{(6)}_i
               + \frac{h^6}{720}f^{(7)}_i + \Ord{h^7} \\
  f''_{i\pm1} &= f''_i \pm hf^{(3)}_i + \frac{h^2}{2}f^{(4)}_i \pm \frac{h^3}{6}f^{(5)}_i
                + \frac{h^4}{24}f^{(6)}_i \pm \frac{h^5}{120}f^{(7)}_i + \Ord{h^6}
\end{align}
関数値については：
\begin{align}
  f_{i+1} - f_{i-1} &= 2hf'_i + \frac{h^3}{3}f^{(3)}_i + \frac{h^5}{60}f^{(5)}_i
                       + \frac{h^7}{2520}f^{(7)}_i + \Ord{h^8}
\end{align}

\subsubsection*{左辺（LHS）の展開}

\begin{align}
  \text{LHS} &= \alpha_1 f'_{i-1} + f'_i + \alpha_1 f'_{i+1}
  = (1 + 2\alpha_1)f'_i + \alpha_1 h^2 f^{(3)}_i
    + \frac{\alpha_1}{12}h^4 f^{(5)}_i + \frac{\alpha_1}{360}h^6 f^{(7)}_i + \Ord{h^7}
\end{align}

\subsubsection*{右辺（RHS）の展開}

\begin{align}
  \frac{a_1}{h}(f_{i+1}-f_{i-1})
  &= 2a_1 f'_i + \frac{a_1}{3}h^2 f^{(3)}_i + \frac{a_1}{60}h^4 f^{(5)}_i
     + \frac{a_1}{2520}h^6 f^{(7)}_i + \Ord{h^7} \\
  b_1 h\bigl(f''_{i+1}-f''_{i-1}\bigr)
  &= 2b_1 h^2 f^{(3)}_i + \frac{b_1}{3}h^4 f^{(5)}_i
     + \frac{b_1}{60}h^6 f^{(7)}_i + \Ord{h^7}
\end{align}

合計：
\begin{equation}
  \text{RHS} = 2a_1 f'_i
  + \Bigl(\frac{a_1}{3} + 2b_1\Bigr)h^2 f^{(3)}_i
  + \Bigl(\frac{a_1}{60} + \frac{b_1}{3}\Bigr)h^4 f^{(5)}_i
  + \Bigl(\frac{a_1}{2520} + \frac{b_1}{60}\Bigr)h^6 f^{(7)}_i + \Ord{h^7}
\end{equation}

\subsubsection*{係数比較と連立方程式}

$\Ord{h^6}$ まで LHS $=$ RHS が成立するには，
$f'_i$，$h^2 f^{(3)}_i$，$h^4 f^{(5)}_i$ の各係数が一致しなければならない：

\begin{align}
  f'_i       &:\quad 1 + 2\alpha_1 = 2a_1 \label{eq:app_I_eq1} \\
  h^2 f^{(3)}_i &:\quad \alpha_1 = \frac{a_1}{3} + 2b_1 \label{eq:app_I_eq2} \\
  h^4 f^{(5)}_i &:\quad \frac{\alpha_1}{12} = \frac{a_1}{60} + \frac{b_1}{3} \label{eq:app_I_eq3}
\end{align}

\subsubsection*{解法}

式~\eqref{eq:app_I_eq1} より $a_1 = (1+2\alpha_1)/2$．
式~\eqref{eq:app_I_eq2} より $b_1 = (\alpha_1 - a_1/3)/2$：
\[
  b_1 = \frac{1}{2}\!\left(\alpha_1 - \frac{1+2\alpha_1}{6}\right)
      = \frac{1}{2}\cdot\frac{6\alpha_1 - 1 - 2\alpha_1}{6}
      = \frac{4\alpha_1 - 1}{12}
\]
$a_1$ と $b_1$ を式~\eqref{eq:app_I_eq3} に代入：
\[
  \frac{\alpha_1}{12}
  = \frac{1+2\alpha_1}{120} + \frac{4\alpha_1-1}{36}
\]
両辺に $360$ を掛けて整理：
\[
  30\alpha_1 = 3(1+2\alpha_1) + 10(4\alpha_1-1)
             = 3 + 6\alpha_1 + 40\alpha_1 - 10
             = 46\alpha_1 - 7
\]
\[
  -16\alpha_1 = -7
  \quad\Longrightarrow\quad
  \alpha_1 = \frac{7}{16}
\]

以上より：
\begin{alignat}{3}
  \alpha_1 &= \frac{7}{16}, \qquad
  &a_1 &= \frac{1+2\cdot(7/16)}{2} = \frac{1 + 7/8}{2} = \frac{15}{16}, \qquad
  &b_1 &= \frac{4\cdot(7/16)-1}{12} = \frac{7/4-1}{12} = \frac{3/4}{12} = \frac{1}{16}
\end{alignat}

\subsubsection*{打ち切り誤差}

$h^6 f^{(7)}_i$ の LHS 係数 $-$ RHS 係数：
\[
  TE_{\mathrm{I}} = \left(\frac{\alpha_1}{360} - \frac{a_1}{2520} - \frac{b_1}{60}\right)h^6 f^{(7)}_i
\]
数値代入：
\begin{align*}
  &\frac{7/16}{360} - \frac{15/16}{2520} - \frac{1/16}{60}
  = \frac{1}{16}\!\left(\frac{7}{360} - \frac{15}{2520} - \frac{1}{60}\right) \\
  &= \frac{1}{16}\!\left(\frac{49}{2520} - \frac{15}{2520} - \frac{42}{2520}\right)
  = \frac{1}{16}\cdot\frac{-8}{2520}
  = -\frac{1}{5040} = -\frac{1}{7!}
\end{align*}
したがって
\[
  TE_{\mathrm{I}} = -\frac{1}{7!}\,h^6 f^{(7)}_i = \Ord{h^6}
\]
が成立し，スキームが $\Ord{h^6}$ 精度であることが確認された．\qed

\subsubsection{Equation-II の係数導出}
\label{app:ccd_coef_derivation_II}

Equation-II（式~\eqref{eq:CCD_II}）の各項を展開する．
二次微分値と関数値の差分については：
\begin{align}
  f''_{i\pm1} &= f''_i \pm hf^{(3)}_i + \frac{h^2}{2}f^{(4)}_i \pm \frac{h^3}{6}f^{(5)}_i
                + \frac{h^4}{24}f^{(6)}_i \pm \frac{h^5}{120}f^{(7)}_i
                + \frac{h^6}{720}f^{(8)}_i + \Ord{h^7} \\
  f_{i-1} - 2f_i + f_{i+1} &= h^2 f''_i + \frac{h^4}{12}f^{(4)}_i
                               + \frac{h^6}{360}f^{(6)}_i
                               + \frac{h^8}{20160}f^{(8)}_i + \Ord{h^8} \\
  f'_{i+1} - f'_{i-1} &= 2hf''_i + \frac{h^3}{3}f^{(4)}_i
                         + \frac{h^5}{60}f^{(6)}_i
                         + \frac{h^7}{2520}f^{(8)}_i + \Ord{h^8}
\end{align}

\subsubsection*{左辺（LHS）の展開}

\begin{equation}
  \text{LHS} = \beta_2 f''_{i-1} + f''_i + \beta_2 f''_{i+1}
  = (1+2\beta_2)f''_i + \beta_2 h^2 f^{(4)}_i
    + \frac{\beta_2}{12}h^4 f^{(6)}_i + \frac{\beta_2}{360}h^6 f^{(8)}_i + \Ord{h^7}
\end{equation}

\subsubsection*{右辺（RHS）の展開}

\begin{align}
  \frac{a_2}{h^2}(f_{i-1}-2f_i+f_{i+1})
  &= a_2 f''_i + \frac{a_2}{12}h^2 f^{(4)}_i + \frac{a_2}{360}h^4 f^{(6)}_i
     + \frac{a_2}{20160}h^6 f^{(8)}_i + \Ord{h^7} \\
  \frac{b_2}{h}(f'_{i+1}-f'_{i-1})
  &= 2b_2 f''_i + \frac{b_2}{3}h^2 f^{(4)}_i + \frac{b_2}{60}h^4 f^{(6)}_i
     + \frac{b_2}{2520}h^6 f^{(8)}_i + \Ord{h^7}
\end{align}

合計：
\begin{equation}
  \text{RHS} = (a_2+2b_2)f''_i
  + \Bigl(\frac{a_2}{12}+\frac{b_2}{3}\Bigr)h^2 f^{(4)}_i
  + \Bigl(\frac{a_2}{360}+\frac{b_2}{60}\Bigr)h^4 f^{(6)}_i
  + \Bigl(\frac{a_2}{20160}+\frac{b_2}{2520}\Bigr)h^6 f^{(8)}_i + \Ord{h^7}
\end{equation}

\subsubsection*{係数比較と連立方程式}

$f''_i$，$h^2 f^{(4)}_i$，$h^4 f^{(6)}_i$ の係数を比較する：
\begin{align}
  f''_i        &:\quad (1+2\beta_2) - (a_2+2b_2) = 0 \label{eq:app_II_eq1} \\
  h^2 f^{(4)}_i &:\quad \beta_2 - \Bigl(\frac{a_2}{12}+\frac{b_2}{3}\Bigr) = 0 \label{eq:app_II_eq2} \\
  h^4 f^{(6)}_i &:\quad \frac{\beta_2}{12} - \Bigl(\frac{a_2}{360}+\frac{b_2}{60}\Bigr) = 0 \label{eq:app_II_eq3}
\end{align}

\subsubsection*{解法}

式~\eqref{eq:app_II_eq2} および \eqref{eq:app_II_eq3} から $\beta_2$ を消去する．
式~\eqref{eq:app_II_eq2} の12倍 $=$ 式~\eqref{eq:app_II_eq3} の左辺に代入：
\[
  \frac{a_2}{12} + \frac{b_2}{3} = \frac{a_2}{30} + \frac{b_2}{5}
\]
整理：
\[
  a_2\!\left(\frac{1}{12}-\frac{1}{30}\right) = b_2\!\left(\frac{1}{5}-\frac{1}{3}\right)
  \quad\Longrightarrow\quad
  a_2\cdot\frac{3}{60} = b_2\cdot\!\left(-\frac{2}{15}\right)
  \quad\Longrightarrow\quad
  b_2 = -\frac{3a_2}{8}
\]
式~\eqref{eq:app_II_eq1} に代入：
\[
  1 + 2\beta_2 = a_2 + 2b_2 = a_2 - \frac{3a_2}{4} = \frac{a_2}{4}
  \quad\Longrightarrow\quad
  \beta_2 = \frac{a_2}{8} - \frac{1}{2}
\]
これを式~\eqref{eq:app_II_eq2} に代入：
\[
  \frac{a_2}{8} - \frac{1}{2} = \frac{a_2}{12} + \frac{1}{3}\cdot\!\left(-\frac{3a_2}{8}\right)
  = \frac{a_2}{12} - \frac{a_2}{8}
  = -\frac{a_2}{24}
\]
\[
  \frac{a_2}{8} + \frac{a_2}{24} = \frac{1}{2}
  \quad\Longrightarrow\quad
  \frac{4a_2}{24} = \frac{1}{2}
  \quad\Longrightarrow\quad
  a_2 = 3
\]

以上より：
\begin{alignat}{3}
  a_2 &= 3, \qquad
  &b_2 &= -\frac{3\times3}{8} = -\frac{9}{8}, \qquad
  &\beta_2 &= \frac{3}{8} - \frac{1}{2} = -\frac{1}{8}
\end{alignat}

\subsubsection*{打ち切り誤差}

$h^6 f^{(8)}_i$ の LHS 係数 $-$ RHS 係数：
\[
  TE_{\mathrm{II}} = \left(\frac{\beta_2}{360} - \frac{a_2}{20160} - \frac{b_2}{2520}\right)h^6 f^{(8)}_i
\]
数値代入：
\[
  \frac{-1/8}{360} - \frac{3}{20160} - \frac{-9/8}{2520}
  = -\frac{1}{2880} - \frac{3}{20160} + \frac{9}{20160}
  = -\frac{7}{20160} + \frac{6}{20160}
  = -\frac{1}{20160}
\]
したがって
\[
  TE_{\mathrm{II}} = -\frac{1}{20160}\,h^6 f^{(8)}_i = \Ord{h^6}
\]
が成立し，Equation-II も $\Ord{h^6}$ 精度であることが確認された．\qed

\medskip
\noindent
主要打ち切り誤差係数 $-1/5040 = -1/7!$（Eq-I）と $-1/20160$（Eq-II）はいずれも
7点陽的スキームの打ち切り誤差係数と同オーダーであり，
3点ステンシルで7点相当の精度を実現していることを示している．

% ═══════════════════════════════════════════════════════════════════════════════
